Retrieval and adaptation address different parts of an evidence workflow. Retrieval-augmented generation combines a retriever with a generative model, while retrieval can also be incorporated into language-model training \citep{lewis2020rag,guu2020realm}. Accordingly, a lexical ranking or nearest-neighbor answer copy is a useful baseline, but is not automatically a generative RAG result. LoRA changes trainable low-rank matrices while freezing base weights \citep{hu2022lora}; a changed prompt or retrieval index is not such adaptation. Our comparisons therefore identify the actual executed arm, source access, checkpoint and decoding conditions.

Prior comparisons do not establish a universal winner for this corpus. \citet{ovadia2024finetuning} compared unsupervised fine-tuning with retrieval under a particular factual-knowledge design, while \citet{gekhman2024hallucinations} studied known and unknown factual examples and training duration in controlled closed-book question answering. Their results motivate controlled baselines and development-set selection, rather than assuming that fitting channeling assertions supplies factual knowledge or that any RAG system must outperform supervised adaptation.

Evaluation must separate a relevant retrieved context, faithfulness to that context and relevance to the question. RAGAs operationalizes these dimensions through model-assisted metrics \citep{es2024ragas}; ARES uses independent human annotations to calibrate trained judges and prediction-powered estimates \citep{saadfalcon2024ares}. Without corresponding human labels, our automatic overlap and retrieval measurements remain proxies. They do not measure the truth of source claims or inherit calibrated human reliability. Context-position sensitivity \citep{liu2024lostmiddle} and train--test duplication \citep{lee2022dedup} further motivate explicit evidence windows, token budgets and source-family leakage audits.
